Lengthening fire-weather seasons and shifting vegetation flammability are driving unprecedented fire activity across Mediterranean Europe and other fire-prone regions. Portugal exemplifies this trend: extreme fire-weather indices have risen steadily since the late twentieth century, and projections indicate the most severe fire-weather class will cover much of the country by mid-century. Effective fire management demands rapid detection at scales fine enough to resolve individual fire fronts, yet no current or planned satellite system delivers both sub-hour revisit and metre-class resolution where it is needed most.

Geostationary sensors scan continental areas every 10--15~minutes but resolve only 2--3~km per pixel, while low-Earth-orbit missions achieve sub-30~m resolution at the cost of multi-day revisit intervals. This resolution--revisit--latency gap leaves fire-prone territories without timely, high-resolution space-based fire intelligence. The present work asks: \emph{What constellation design can close this gap, achieving roughly 30-minute revisit and approximately 6~m ground resolution over continental Portugal?} A secondary question is: \emph{What coverage does the constellation provide over global fire-prone regions?}

The study integrates fire-remote-sensing science, orbital mechanics, and high-fidelity numerical simulation.
A core study domain was first delineated by to Portugal's territory according it's fire season pattern, than a global study domain was defined to current fire-weather severity, establishing the geographical, spatial, and temporal requirements the constellation must satisfy.
An analytical survey of repeating ground-track orbits (RGTOs), generalised to arbitrary target latitudes, produced a family of RGTO modes for daily revisit in continental Portugal; from this group of orbital geometries the ratio $\tau = 15/1$ at approximately 488~km altitude and inclination ${\sim}\,40.2°$~was selected as best matched to the target latitude core study domain at center of Portugal, an resonating orbit around the city of Coimbra, providing the fastest revisit possible.
This selected orbitalmode was embedded in a Walker constellation framework. Parametric trades then explored the number of orbital planes, satellites per plane, and maximum sensor off-nadir viewing angle within a Walker-delta framework were allowed to optmized the revisit over the Portuguese core study domain.
The candidate constellations were propagated sing the high-fidelity General Mission Analysis Tool (GMAT) --- with a $8\times8$ geopotential, lunisolar perturbations, and atmospheric drag--- validating the orbital stability, observational-link pattern, temporal coverage, revisit statistics, and ground sampling distance over both the Portuguese core domain and a global fire-prone domain. 

The selected constellation---designated $C_{\text{Zoom}} = C_{15/5/72}$, comprising 15~satellites in 5~equally spaced orbital planes phased by $72^\circ$ of right ascending node, with 3 satellite per plane distributed $120^\circ$ of each by shifts in the true anomaly---achieves a mean revisit of approximately 21~minutes over Portugal with ground resolution below 6~m per pixel.
The repeating ground-track design delivers roughly $1.75\times$ the daily coverage time of a conventional sun-synchronous orbit at comparable altitude, effectively halving the fleet size that would otherwise be required.
A sensitivity analysis confirms that a $65°$ off-nadir sensor limit is essential: tightening it to cover only Portugal $30°$ keeping same cadence nearly doubles the required fleet to 27~satellites, underscoring wide-angle pointing as a key design driver for keeping the constellation compact.
Designed for Portugal, the coverage footprint extends naturally to all major fire-prone regions with exception of boreal region: the temperate belt---California, the broader Mediterranean basin, southern Africa, and south-eastern Australia---receives 6--7~hours of daily observation with worst-case revisit gaps of about 30 minutes, while tropical regions such as Amazonia, central Africa, and south-east Asia retain at least 3.5~hours per day with an average revisit around 40 to 50 minutes.
To expand the narrow field of view that metre-class resolution demands on a small platform, a \emph{Watcher--Zoom} tip-and-cue architecture is proposed: five existing geostationary fire-detection satellites provide coarse alerts every 10--15~minutes, cueing the nearest low-orbit spacecraft to steer its imager toward the hotspot and acquire high-resolution imagery.

No surveyed fire-detection constellation combines sub-30-minute revisit with metre-scale resolution over any examined fire-prone region. The architecture proposed here---fifteen satellites on a resonant orbit, guided by geostationary fire-watch assets---closes this gap for Portugal and extends meaningful coverage to every major fire-prone region on Earth, turning a national mission into a global cooperative-observation platform.


The study integrates fire-remote-sensing science, orbital mechanics, and high-fidelity numerical simulation. A core study domain was delineated over Portugal according to its fire-season pattern; a global domain was then defined by current fire-weather severity, establishing the geographical, spatial, and temporal requirements. An analytical survey of repeating ground-track orbits (RGTOs), generalised to arbitrary target latitudes, produced a family of modes for daily revisit over Portugal; from this set the ratio $\tau = 15/1$ at approximately 488~km altitude and ${\sim}\,40.2°$ inclination was selected as best matched to the core domain centred on Coimbra, providing the fastest revisit. This orbital mode was embedded in a Walker-delta framework; parametric trades explored the number of planes, satellites per plane, and maximum off-nadir angle to optimise revisit. Candidate constellations were propagated with the General Mission Analysis Tool (GMAT)---$8\times8$ geopotential, lunisolar perturbations, and atmospheric drag---validating orbital stability, coverage, revisit statistics, and ground sampling distance over both domains.

The selected constellation---$C_{\text{Zoom}} = C_{15/5/72}$, 15~satellites in 5~planes phased by $72^\circ$, 3~per plane at $120^\circ$ true-anomaly spacing---achieves a mean revisit of approximately 21~minutes over Portugal with ground resolution below 6~m. The repeating ground-track delivers roughly $1.75\times$ the daily coverage of a sun-synchronous orbit at comparable altitude, effectively halving the required fleet. A sensitivity analysis confirms that the $65°$ off-nadir limit is essential: tightening it to $30°$ nearly doubles the fleet to 27~satellites, underscoring wide-angle pointing as a key design driver. The coverage extends naturally to all major fire-prone regions except the boreal belt: temperate zones---California, the broader Mediterranean, southern Africa, south-eastern Australia---receive 6--7~hours of daily observation with worst-case gaps of about 30~minutes, while tropical regions retain at least 3.5~hours per day with average revisit of 40--50~minutes. To compensate the narrow field of view that metre-class resolution demands, a \emph{Watcher--Zoom} tip-and-cue architecture is proposed: five geostationary fire-detection satellites provide alerts every 10--15~minutes, cueing the nearest low-orbit spacecraft to image the hotspot at high resolution.

% \vspace{1em}
% \noindent\textbf{Keywords:} satellite constellation, fire detection, repeating ground-track orbit, mission design, fire detection, Portugal, Earth observation